\documentclass[sigconf]{acmart}
\citestyle{acmauthoryear}
\setcopyright{none}

\title{UNIX V4 and the Utah Teapot}
\author{Thalia Archibald}
\affiliation{%
  \institution{University of Utah}
  \department{Kahlert School of Computing}
  \city{Salt Lake City}
  \state{Utah}
  \country{USA}
}
\email{thalia.archibald@utah.edu}
\renewcommand{\shortauthors}{Archibald}

\begin{document}
\begin{abstract}
We recently recovered UNIX V4 from a 1974 magnetic tape at the University of
Utah. This version of the UNIX operating system, thought to have been lost, was
the 19th copy distributed to the public, just months after the first public
announcement. It was originally acquired by Martin Newell while managing the
computer graphics laboratory, and it was likely connected to his foundational
research in procedural modeling and the famous Utah teapot. This talk shows the
unique insight this artifact provides into the dawn of computer graphics and the
history of UNIX.
\end{abstract}

\maketitle
\section{Tape Recovery}

\section{UNIX History}

UNIX is an influential time-sharing operating system created by Ken Thompson,
Dennis Ritchie, and others at Bell Laboratories in the 1970s.

\section{Research at the University of Utah}

That phone number is still live today, used as the Kahlert School of Computing
front desk number.

Months prior, a professor from the Department of Chemistry inquired about UNIX,
to support pioneering research in mass spectroscopy.


\end{document}
